% Frozen candidate addendum: reduction invariance and the nonreduced
% extension of P21.3 curve rigidity.
% Date: 27 August 2026.
% This packet is not integrated into the live book.

\begin{proposition}
\label{prop:p21p3-curve-rigidity-nonreduced}
Let \(K\) be an algebraically closed complete nontrivially valued
non-Archimedean field and let \(X\) be a good separated strict
\(K\)-analytic space, pure of dimension one and without boundary.  The
space \(X\) is not assumed reduced.  Write

\[
             \iota:X_{\mathrm{red}}\hookrightarrow X
\]

for the canonical nilpotent closed immersion.  Let \((X_i)_i\) be the
irreducible components of \(X_{\mathrm{red}}\), assume that their
normalizations are rig-smooth, and write

\[
 \nu_i:\widetilde X_i\longrightarrow X_i
       \hookrightarrow X_{\mathrm{red}}\xrightarrow{\iota}X,
 \qquad
 \delta_i:=(\nu_i)_*\delta_{\widetilde X_i}.
\]

Let \(m_i\in\mathbf Z_{>0}\) be the coefficient of \([X_i]\) in the
fundamental cycle \([X]\).

Then the following assertions hold.

\begin{enumerate}
\item\label{itm:p21p3-nr:calculus}
For every \((p,q)\), pullback by \(\iota\) identifies the sheaves of
ordinary CLD-smooth forms on \(X\) and \(X_{\mathrm{red}}\), compatibly
with wedge product, \(J\), \(\mathrm d'\), \(\mathrm d''\), support and all
three positivity cones.  On compactly supported test forms this is an
isomorphism of topological vector spaces.  Its transpose identifies the
CLD current sheaves, their weak topologies, positivity, and
\(\mathrm d'\)- and \(\mathrm d''\)-closedness.  It identifies
multiplication by smooth functions for all currents.  If \(S\) is positive
and \(f\) is a locally bounded Borel function, then it also identifies the
CLD product \(fS\); by linearity the same holds for currents locally
expressible as differences of positive currents.  No multiplication by a
Borel function is asserted for an arbitrary order-zero current.  Under the
common underlying topological space, the CLD-psh and continuous CLD-psh
sheaves on \(X\) and \(X_{\mathrm{red}}\) coincide.  The analogous pullback
identification holds for \(G\)-forms.

\item\label{itm:p21p3-nr:integration}
The integration current is not, in general, invariant under reduction.
Precisely,

\begin{equation}
\label{eq:p21p3-nr-fundamental-cycle}
        \delta_X=\sum_i m_i\delta_i,
        \qquad
        \iota_*\delta_{X_{\mathrm{red}}}=\sum_i\delta_i.
\end{equation}

Both sums are locally finite.  In particular,
\(\delta_X=\iota_*\delta_{X_{\mathrm{red}}}\) if and only if \(X\) is
generically reduced along every \(X_i\).

\item\label{itm:p21p3-nr:classification}
Every positive \(\mathrm d'\)- and \(\mathrm d''\)-closed CLD current \(T\)
of bidegree \((0,0)\) on \(X\) has a unique expression

\begin{equation}
\label{eq:p21p3-nr-classification}
                     T=\sum_i c_i\delta_i,
                     \qquad c_i\in\mathbf R_{\geq0}.
\end{equation}

The sum is locally finite.  If \(X_{\mathrm{red}}\) is connected and
rig-smooth, there is one generic multiplicity \(m\), and every such current
has the form \(T=c\delta_X\) for a unique \(c\geq0\).

\item\label{itm:p21p3-nr:bedford-taylor}
Let \(W\subseteq X\) be open and let \(U,V\) be continuous CLD-psh
functions on \(W\).  Identify \(W\) with its inverse image in
\(X_{\mathrm{red}}\), and, for every continuous function \(F\) on \(W\),
put
\[
 [F]_X:=F(\delta_X|_W),\qquad
 [F]_{X_{\mathrm{red}}}:=F(\delta_{X_{\mathrm{red}}}|_W).
\]
With the identifications in \cref{itm:p21p3-nr:calculus}, one has the exact
equivalence

\begin{equation}
\label{eq:p21p3-nr-curvature-equivalence}
 \mathrm d'\mathrm d''[U]_X=\mathrm d'\mathrm d''[V]_X
 \quad\Longleftrightarrow\quad
 \mathrm d'\mathrm d''[U]_{X_{\mathrm{red}}}
   =\mathrm d'\mathrm d''[V]_{X_{\mathrm{red}}}.
\end{equation}

If these equivalent conditions hold, then every \(T\) in
\cref{itm:p21p3-nr:classification} satisfies

\begin{equation}
\label{eq:p21p3-nr-bt-conclusion}
       \mathrm d'\mathrm d''\bigl((U-V)T|_W\bigr)=0.
\end{equation}

Consequently every positive closed CLD current of bidegree \((0,0)\) on
\(X\) is a Bedford--Taylor current.  Equivalently, under the current
identification of \cref{itm:p21p3-nr:calculus}, the Bedford--Taylor
property is invariant under the nilpotent thickening
\(X_{\mathrm{red}}\hookrightarrow X\) in this one-dimensional setting.
\end{enumerate}
\end{proposition}

\begin{remark}
\label{rem:p21p3-nr-double-line}
For

\[
 X=\bigl(\operatorname{Spec}K[T,\epsilon]/(\epsilon^2)\bigr)^{\mathrm{an}},
 \qquad X_{\mathrm{red}}=\mathbf A_K^{1,\mathrm{an}},
\]

the unique generic multiplicity is \(2\).  Hence all smooth forms and all
positive closed currents are identified with those on the analytic line,
whereas
\(\delta_X=2\iota_*\delta_{\mathbf A_K^{1,\mathrm{an}}}\).
Thus generic multiplicity, and not any new nilpotent differential form or
current, is the precise obstruction.  The proof of
\cref{prop:p21p3-curve-rigidity-nonreduced} uses only the normalizations of
the components of \(X_{\mathrm{red}}\); no normalization of the
nonreduced space is needed.
\end{remark}
